% GENERATED FILE. Edit canonical lessons, then run npm run generate:books.
% canonical-source-hash: fnv1a64:dcaa3df792784dc0
% canonical-lessons: MW-C145-gulabi, MW-C145-kesariyo, MW-C145-asmani, MW-C145-baingni, MW-C145-narangi

\chapter{Pink, Saffron-coloured, Sky-blue, Purple, Orange}
\label{ch:mw-pink-saffron-coloured-sky-blue-purple-orange}
\hlchaptermodality{145}

\begin{chapteropening}
\textbf{By the end of this chapter:} \emph{I can say pink, saffron-coloured, sky-blue, purple and orange.}

Complete the last lesson of chapter 145: I can say pink, saffron-coloured, sky-blue, purple and orange.
\end{chapteropening}

\section[\texorpdfstring{\mw{गुलाबी} (gulābī)}{गुलाबी (gulābī)}]{\mw{गुलाबी} (gulābī) — pink}

\label{lesson:MW-C145-gulabi}



\begin{quote}
Before the new one: say the Marwadi for thin, then the Marwadi for short.
\end{quote}

\subsection*{You'll want to know: \mw{गुलाबी}}
\textbf{\mw{गुलाबी}} — \emph{gulābī} — \textquotedblleft{}pink\textquotedblright{}.

\begin{cousinweb}
\textbf{\mw{गुलाबी}} (\emph{gulābī}) is a loan, usually traced to Persian \emph{gulābī}, \textquotedblleft{}rose-coloured\textquotedblright{}. It is built on Persian \emph{gulāb}, \textquotedblleft{}rose-water\textquotedblright{}.
\end{cousinweb}

\begin{grammarlens}[title={What you've built}]
One more word for this chapter.
\end{grammarlens}

\subsection*{Your turn}
\begin{itemize}
\raggedright
  \item \emph{Say it:} \emph{gulābī}
  \item \emph{Say it:} \emph{gulābī}, once more
  \item \emph{Say it:} say \emph{gulābī}
  \item \emph{Recall:} say the Marwadi for thin, then the Marwadi for short, then say \emph{gulābī} again
\end{itemize}

\begin{tcolorbox}[breakable,colback=teal!4,colframe=teal!35!black,title={Before you move on}]
What does \mw{गुलाबी} mean? (\textquotedblleft{}Pink\textquotedblright{}.) Say it once more.
\end{tcolorbox}
\section[\texorpdfstring{\mw{केसरियो} (kesariyo)}{केसरियो (kesariyo)}]{\mw{केसरियो} (kesariyo) — saffron-coloured}

\label{lesson:MW-C145-kesariyo}



\begin{quote}
Before the new one: say the Marwadi for short, then the Marwadi for pink.
\end{quote}

\subsection*{You'll want to know: \mw{केसरियो}}
\textbf{\mw{केसरियो}} — \emph{kesariyo} — \textquotedblleft{}saffron-coloured\textquotedblright{}.

\begin{cousinweb}
\textbf{\mw{केसरियो}} (\emph{kesariyo}) is built with the suffix \emph{-iyo} on \emph{kesar}, \textquotedblleft{}saffron\textquotedblright{}, which CDIAL lists as Gujarati and Hindi \emph{kesar} under Sanskrit \emph{kesara}, \textquotedblleft{}hair, filament\textquotedblright{} (CDIAL 3474): saffron is the flower's filament.
\end{cousinweb}

\begin{grammarlens}[title={What you've built}]
One more word for this chapter.
\end{grammarlens}

\subsection*{Your turn}
\begin{itemize}
\raggedright
  \item \emph{Say it:} \emph{kesariyo}
  \item \emph{Say it:} \emph{kesariyo}, once more
  \item \emph{Say it:} say \emph{kesariyo}
  \item \emph{Recall:} say the Marwadi for short, then the Marwadi for pink, then say \emph{kesariyo} again
\end{itemize}

\begin{tcolorbox}[breakable,colback=teal!4,colframe=teal!35!black,title={Before you move on}]
What does \mw{केसरियो} mean? (\textquotedblleft{}Saffron-coloured\textquotedblright{}.) Say it once more.
\end{tcolorbox}
\section[\texorpdfstring{\mw{आसमानी} (āsmānī)}{आसमानी (āsmānī)}]{\mw{आसमानी} (āsmānī) — sky-blue}

\label{lesson:MW-C145-asmani}



\begin{quote}
Before the new one: say the Marwadi for pink, then the Marwadi for saffron-coloured.
\end{quote}

\subsection*{You'll want to know: \mw{आसमानी}}
\textbf{\mw{आसमानी}} — \emph{āsmānī} — \textquotedblleft{}sky-blue\textquotedblright{}.

\begin{cousinweb}
\textbf{\mw{आसमानी}} (\emph{āsmānī}) is a loan, usually traced to Persian \emph{āsmānī}, \textquotedblleft{}sky-coloured\textquotedblright{}. It is built on Persian \emph{āsmān}, \textquotedblleft{}sky\textquotedblright{}.
\end{cousinweb}

\begin{grammarlens}[title={What you've built}]
One more word for this chapter.
\end{grammarlens}

\subsection*{Your turn}
\begin{itemize}
\raggedright
  \item \emph{Say it:} \emph{āsmānī}
  \item \emph{Say it:} \emph{āsmānī}, once more
  \item \emph{Say it:} say \emph{āsmānī}
  \item \emph{Recall:} say the Marwadi for pink, then the Marwadi for saffron-coloured, then say \emph{āsmānī} again
\end{itemize}

\begin{tcolorbox}[breakable,colback=teal!4,colframe=teal!35!black,title={Before you move on}]
What does \mw{आसमानी} mean? (\textquotedblleft{}Sky-blue\textquotedblright{}.) Say it once more.
\end{tcolorbox}
\section[\texorpdfstring{\mw{बैंगणी} (baiṅgṇī)}{बैंगणी (baiṅgṇī)}]{\mw{बैंगणी} (baiṅgṇī) — purple}

\label{lesson:MW-C145-baingni}



\begin{quote}
Before the new one: say the Marwadi for saffron-coloured, then the Marwadi for sky-blue.
\end{quote}

\subsection*{You'll want to know: \mw{बैंगणी}}
\textbf{\mw{बैंगणी}} — \emph{baiṅgṇī} — \textquotedblleft{}purple\textquotedblright{}.

\begin{cousinweb}
\textbf{\mw{बैंगणी}} (\emph{baiṅgṇī}) is built with \emph{-ī} on \emph{baiṅgaṇ}, \textquotedblleft{}aubergine\textquotedblright{}, from Sanskrit \emph{vātiṅgaṇa}, which CDIAL thinks may come from Dravidian (CDIAL 11503). CDIAL lists Hindi \emph{baĩjnī}, \textquotedblleft{}purple\textquotedblright{}, under the same word.
\end{cousinweb}

\begin{grammarlens}[title={What you've built}]
One more word for this chapter.
\end{grammarlens}

\subsection*{Your turn}
\begin{itemize}
\raggedright
  \item \emph{Say it:} \emph{baiṅgṇī}
  \item \emph{Say it:} \emph{baiṅgṇī}, once more
  \item \emph{Say it:} say \emph{baiṅgṇī}
  \item \emph{Recall:} say the Marwadi for saffron-coloured, then the Marwadi for sky-blue, then say \emph{baiṅgṇī} again
\end{itemize}

\begin{tcolorbox}[breakable,colback=teal!4,colframe=teal!35!black,title={Before you move on}]
What does \mw{बैंगणी} mean? (\textquotedblleft{}Purple\textquotedblright{}.) Say it once more.
\end{tcolorbox}
\section[\texorpdfstring{\mw{नारंगी} (nāraṅgī)}{नारंगी (nāraṅgī)}]{\mw{नारंगी} (nāraṅgī) — orange}

\label{lesson:MW-C145-narangi}



\begin{quote}
Before the new one: say the Marwadi for sky-blue, then the Marwadi for purple.
\end{quote}

\subsection*{You'll want to know: \mw{नारंगी}}
\textbf{\mw{नारंगी}} — \emph{nāraṅgī} — \textquotedblleft{}orange\textquotedblright{}.

\begin{cousinweb}
\textbf{\mw{नारंगी}} (\emph{nāraṅgī}) is a tadbhava, a word inherited from Sanskrit and reshaped by sound change. It comes from \emph{nāraṅga}, \textquotedblleft{}orange tree\textquotedblright{}. CDIAL gives no Marwari form in this entry (CDIAL 7073); the match is through Hindi \emph{nāraṅgī} and Gujarati \emph{nārãgī}. The colour is named for the fruit.
\end{cousinweb}

\begin{grammarlens}[title={What you've built}]
One more word for this chapter.
\end{grammarlens}

\subsection*{Your turn}
\begin{itemize}
\raggedright
  \item \emph{Say it:} \emph{nāraṅgī}
  \item \emph{Say it:} \emph{nāraṅgī}, once more
  \item \emph{Say it:} say \emph{nāraṅgī}
  \item \emph{Recall:} say the Marwadi for sky-blue, then the Marwadi for purple, then say \emph{nāraṅgī} again
\end{itemize}

\begin{tcolorbox}[breakable,colback=teal!4,colframe=teal!35!black,title={Before you move on}]
What does \mw{नारंगी} mean? (\textquotedblleft{}Orange\textquotedblright{}.) Say it once more.
\end{tcolorbox}
